% migrate-to: execution-market-making, market-making-lob, hft-market-design, market-structure-execution-quality
\subsection{Volatility and Liquidity Interaction}

% migrated: Volatility Is a Liquidity Price -> res:spread-determinants (Practice, Pitfalls); Volatility Source Segmentation, Fundamental Volatility Driver Segmentation, Transitory Volatility and Reversal, Volatility Component Decomposition -> def:volatility-components; Bubbles and Crashes, Bubble-Crash Dynamics Segmentation, Price Accelerator -> mech:bubble-crash-dynamics; Crash Attribution -> mech:bubble-crash-dynamics (Pitfalls); Regulating Volatility -> def:volatility-components (Practice); Garman-Klass Intraday Volatility -> meth:garman-klass-volatility; Spread-Volatility Regression, Log-Log Spread-Volatility Model -> meth:spread-volatility-regression; Spread-Volatility Identification -> meth:spread-volatility-regression (Pitfalls) (full text in notes/artefacts/market-making-lob/derivations.tex)

% assign: hft-market-design
\begin{mechanism} \hlt{Liquidity Spiral}\\
A large order can consume depth, widen spreads, and raise volatility. Higher volatility makes liquidity providers reduce size or cancel quotes, forcing the next order to trade through less liquidity and amplifying the move. The state variable is confidence in replenishment: once traders doubt the book will refill, everyone demands more price concession at the same time.
\end{mechanism}
% assign: hft-market-design
\begin{casestudy} \hlt{Flash Crash Mechanism}\\
The 2010 Flash Crash illustrates a liquidity spiral: a large automated futures sell program met thin liquidity, high-frequency traders recycled inventory quickly, aggressive flow migrated into cash equities, spreads widened, and some routers or venues paused.\\
The trigger was an order; the crash mechanism was fragile liquidity under linked-market stress.
\end{casestudy}
% assign: hft-market-design
\begin{model} \hlt{Depth-Driven Flash Crash Toy Model}\\
On one side of the book, let $F_t$ be the market-order arrival rate, $\bar{F}$ the normal refill rate, $r_t=F_t-\bar{F}$ the overflow, and $I_t=1/D_t$ inverse depth. A compact depth-impact dynamic is $\frac{dI_t}{I_t}=\sigma_I dW_t+r_tdt$.\\
Here $D_t$ is market depth, $I_t$ is impact, $\sigma_I$ is liquidity-shock volatility, and $W_t$ captures random liquidity shocks.\\
Positive overflow $r_t$ makes inverse depth compound upward, so each additional trade has larger price impact than the last.
\end{model}
% assign: hft-market-design
\begin{mechanism} \hlt{Hot-Potato Amplification}\\
If a fraction $\rho$ of traders follows volume mechanically, initial trading rate $r_0$ can amplify roughly as $r_t\simeq \frac{1-\rho^t}{1-\rho}r_0$.\\
When $\rho$ is high, overflow consumes depth faster than passive liquidity can refill it. The loop is mechanical: trading creates volume, volume targets force more trading, and depth absorbs the loop until it fails.
\end{mechanism}
% assign: hft-market-design
\begin{taxonomy} \hlt{Flash-Crash Control Segmentation}
\begin{enumerate}[label=\roman*.]
\setlength{\itemsep}{0pt}
\item \hlt{rate controls}: cap participation rates or stop volume-following feedback;
\item \hlt{depth controls}: require or incentivise fast refill near the touch;
\item \hlt{price controls}: pause trading after large price moves;
\item \hlt{cross-market controls}: coordinate futures, ETFs, cash equities, and options;
\item \hlt{data controls}: detect overflow, quote withdrawal, and broken references in real time.
\end{enumerate}
\end{taxonomy}
% assign: market-structure-execution-quality
\begin{caveat} \hlt{Systemic Risk Link}\\
Fragmentation is not systemic risk by itself. It becomes dangerous when market participants react to stress in the same direction, when related venues desynchronise, or when liquidity withdrawal in one instrument propagates to hedges, ETFs, futures, and correlated stocks.
\end{caveat}
% assign: hft-market-design
\begin{taxonomy} \hlt{Shock Propagation Channels}
\begin{enumerate}[label=\roman*.]
\setlength{\itemsep}{0pt}
\item \hlt{cross-asset hedging}: futures, ETFs, options, and cash baskets transmit inventory pressure;
\item \hlt{routing feedback}: routers chase fading liquidity or pause when fills become unreliable;
\item \hlt{quote withdrawal}: market makers widen or cancel when volatility exceeds their risk budget;
\item \hlt{reference-price dependency}: pegged and benchmark orders inherit stale or broken references;
\item \hlt{correlated constraints}: margin, risk limits, and kill switches cause simultaneous de-risking.
\end{enumerate}
\end{taxonomy}
% assign: hft-market-design
\begin{empiricalfinding} \hlt{Trade Count as HFT Proxy}\\
Rising trades per day can proxy for higher high-frequency participation when trade sizes shrink and message intensity rises. It is not sufficient alone: the same count can reflect genuine liquidity, excessive slicing, or stress-induced churn.
\end{empiricalfinding}
